\NumberedItem{6}
\HistoricTitle{\textsc{Gauss} an H. G. \textsc{Grassmann}.}
\par\medskip\noindent\hspace*{.26\linewidth}Euer Wohlgeboren
\par\indent geehrtes Schreiben würde ich schon früher mit meinem gehorsamsten Danke für Ihr gütigst übersandtes Werk[\textsuperscript{*})] erwiedert haben, wenn ich nicht gewünscht hätte, über letzteres erst eine Ansicht fassen zu können. Zunächst bin ich lange vom Buchbinder aufgehalten; hernach, in einem Gedränge von andern heterogenen Arbeiten Ihr Buch durchlaufend glaube ich zu bemerken, dass die Tendenzen desselben \textls[120]{theilweise} denjenigen Wegen begegnen, auf denen ich selbst nun seit fast einem halben Jahrhundert gewandelt bin, und wovon freilich nur ein kleiner Theil 1831 in den Comment. der Göttingischen Societät[\textsuperscript{**})] und noch mehr in den Göttingischen Gelehrten Anzeigen (1831, Stück 64)[\textsuperscript{***})] gleichsam im Vorbeigehen erwähnt ist; nemlich die concentrirte Metaphysik der complexen Grössen, während von der unendlichen Fruchtbarkeit dieses Princips für Untersuchungen räumliche Verhältnisse betreffend zwar vielfältig in meinen Vorlesungen gehandelt[\textsuperscript{\ensuremath{\dagger}})], aber Proben davon nur hin und wieder, und, \textls[120]{als solche} nur dem aufmerksamen Auge erkennbar, bei andern Veranlassungen mitgetheilt sind[\textsuperscript{\ensuremath{\dagger\dagger}})]. Indessen scheint
\par\medskip\noindent\rule{.24\linewidth}{.4pt}\par\smallskip
{\small
\noindent\hspace*{\parindent}[\textsuperscript{*}) \textit{Die Wissenschaft der extensiven Grösse oder die Ausdehnungslehre} von \textsc{Hermann Grassmann}, I. Theil, \textit{Die lineale Ausdehnungslehre}. Leipzig 1844; Gesammelte mathem. und physikal. Werke I, 1, Leipzig 1894.]
\par\noindent\hspace*{\parindent}[\textsuperscript{**}) \textit{Theoria Residuorum biquadraticorum, Commentatio secunda}, Werke II, S. 93, siehe besonders S. 102 ff.]
\par\noindent\hspace*{\parindent}[\textsuperscript{***}) Werke II, S. 169, siehe besonders S. 175 ff.]
\par\noindent\hspace*{\parindent}[\textsuperscript{\ensuremath{\dagger}}) Vergl. die Auszüge aus einer Dezember 1839 bis Ostern 1840 gehaltenen Vorlesung über die \textit{Theorie der imaginären Grössen}: Werke VIII, S. 331—334 und 346—347.]
\par\noindent\hspace*{\parindent}[\textsuperscript{\ensuremath{\dagger\dagger}}) Vergl. die bereits erwähnte Stelle, Werke II, S. 175, wo mit ähnlichen Worten auf die Inauguraldissertation von 1799 (Werke III, S. 1) und auf die Preisschrift von 1822 (Werke IV, S. 189) hingewiesen wird.]
\par}
\clearpage\SourceLeaf{444}\ArithmeticHead{437}{R}{GAUSS AN GRASSMANN.}
\noindent dies nur eine partielle und entferntere Ähnlichkeit in der Tendenz zu sein; und ich sehe wohl, dass um den eigentlichen Kern Ihres Werks herauszufinden, es nöthig sein wird, sich erst mit Ihren eigenthümlichen Terminologien zu familiarisiren. Da aber dazu, bei mir, nothwendig eine von andern Beschäftigungen freiere Zeit erforderlich sein wird, so darf ich jetzt nicht länger anstehen, Ihnen meinen ergebensten Dank für die gefällige Übersendung Ihres Werks auszusprechen, dem ich die Versicherung der besonderen Hochachtung beifüge, mit welcher ich beharre
\par\noindent\hspace*{.54\linewidth}Euer Wohlgeboren
\par\noindent\hspace*{.62\linewidth}ergebenster Diener
\par\noindent\hspace*{\parindent}Göttingen, 14. December 1844.\hfill C. F. \textsc{Gauss}.\hspace*{.06\linewidth}
